\section{模型评价、结果位置与局限}
\subsection{结果表和图的固定位置}
下表先固定正文需要的字段和来源位置，数值待实验终态后补入。它不是实验结果表，也不对当前运行作判断。
\begin{table}[H]\centering\caption{实验结果待补清单}
\begin{tabularx}{\linewidth}{lYY}\toprule
交付项 & 必要字段 & 当前状态\\\midrule
问题一平均曲线 & $K=1\ldots5$ 的单核基准、Makespan、平均加速比 & \pending\\
问题二平均曲线 & 场景 B 的 Makespan、额外搬运和加速比 & \pending\\
问题三对照曲线 & 无 L2、有 Cache 及同核数 Cache 加速比 & \pending\\
逐用例附表 & 图 id、核数、方案、Makespan、额外搬运、状态 & \pending\\
Cache 诊断 & 查询字节、命中字节、命中率、DDR/Cache 服务量 & \pending\\
消融与敏感性 & 初始/最终方案、容量、带宽和候选来源 & \pending\\\bottomrule
\end{tabularx}\end{table}

\subsection{评价维度}
算法评价分为四个互相独立的层次。第一是可行性：方案覆盖、依赖无环、同核顺序和容量约束必须通过评估器。第二是目标表现：在相同图、核数、配置和单核定义下比较 Makespan 及加速比。第三是机制解释：结构搬运、Spill、Cache 服务和查询事件能否解释时间变化。第四是计算成本：候选数量、官方完整调用、失败和运行耗时由 manifest 记录。程序通过不等于科学结论成立；人工复核和正式提交另行记录。

\subsection{已知局限}
有限候选搜索只能给出候选集合中的可行方案，不能直接声称全局最优。轻量代理可能漏掉真实 Spill 或 FIFO 机制收益，因此代理排序反例必须保留。官方事件模拟描述题面给定的处理器抽象，不等于真实 NPU 实机测量。Cache 命中率是字节层代理指标，不能单独等同于总时间收益。实验结果完成后，本文还需根据缺失用例、失败状态和科学复核意见收窄结论。

\subsection{当前结论边界}
截至本稿生成时，能够确定的是题面对象、配置语义、模型公式、候选生成规则和结果交付协议；全量主矩阵及其对照尚在运行，因而不能写出平均加速比、Cache 命中率、消融收益或任何“优于基线”的科学结论。实验结束后只从冻结的结果 CSV 和 manifest 读取数字，并同步摘要、正文、图注、结论和附录。
